\documentclass[11pt,a4paper]{article}
\usepackage[T1]{fontenc}
\usepackage[utf8]{inputenc}
\usepackage{lmodern}
\usepackage[margin=2.4cm]{geometry}
\usepackage{booktabs}
\usepackage{amsmath,amssymb}
\usepackage{graphicx}
\usepackage{caption}
\usepackage{longtable}
\usepackage{microtype}
\usepackage[hidelinks]{hyperref}
\usepackage{enumitem}
\usepackage{titlesec}
\titleformat{\section}{\normalfont\large\bfseries}{\thesection}{0.7em}{}
\titleformat{\subsection}{\normalfont\normalsize\bfseries}{\thesubsection}{0.7em}{}
\setlist{nosep,leftmargin=1.4em}
\captionsetup{font=small,labelfont=bf}
\hyphenation{multi-physics}

\title{\bfseries Capability Verification of a Multiphysics Finite-Element\\
Framework Across Nine Physics Domains}
\author{Carlos Manuel Mendez}
\date{8 August 2026}

\begin{document}
\maketitle

\begin{abstract}
\noindent
We report a verification study of the MOOSE finite-element framework covering
twenty-five solver capabilities across nine physics domains. Each capability is
exercised by an independent input model and compared against either a
closed-form solution (21 cases) or a published numerical benchmark (4 cases),
under an acceptance tolerance committed to version control \emph{before} the
solver was run. All twenty-five cases meet their declared tolerance. The median
case consumes 7.3\% of its error budget; the tightest consumes 87\%; six consume
less than 1\%. We describe the protocol governing the study --- pre-declared
tolerances, sourced reference values, and a requirement that reported quantities
be mutually consistent --- and we report a post-hoc audit of the results. That audit found no incorrect verdict, but it did find one check that
could not have failed, seven results whose headline presentation was misleading,
and three results whose numbers are weaker than a bare pass suggests. We state
plainly what this study does not establish: it is verification, not validation,
and it contains no comparison against experiment.
\end{abstract}

\section{Introduction}

A general-purpose multiphysics framework is a large surface of solver
capabilities, and a user adopting one needs to know which parts of that surface
actually work for their problem class. The usual evidence --- a regression test
suite maintained by the developers --- answers a different question. It shows
that behaviour has not changed, not that it is right.

This study takes the complementary position. Each of twenty-five capabilities is
posed as a small, self-contained problem for which an independent answer exists,
and the framework's result is compared against that answer under a tolerance
fixed in advance. The framework under test is
MOOSE~\cite{permann2020}, snapshot \texttt{20-10-27-41583-g2bd11a08a7}. The
capabilities span solid mechanics and fracture, heat transfer, fluid dynamics,
porous flow, phase-field evolution, mass transport, neutronics, multiphysics
coupling, and optimisation and uncertainty quantification, the last including a monolithic
fluid--structure interaction benchmark.

The distinction between \emph{verification} and \emph{validation} is load-bearing
throughout. Verification asks whether the equations are solved correctly;
validation asks whether the equations describe reality. Everything reported here
is verification. Section~\ref{sec:limits} states the consequences of that
restriction rather than leaving them implicit.

The contribution is not that a mature framework passes tests. It is the
protocol, and the finding that the protocol catches things: an audit of the
study's own results, reported in Section~\ref{sec:audit}, found a check that had
been passing for the wrong reason.

\section{Method}

\subsection{Case structure}
Each case is a directory containing the MOOSE input files, a
\texttt{reference.json} declaring the reference value and the acceptance
tolerance, a verifier script, and a \texttt{result.json} written by that
verifier. The summary table in Section~\ref{sec:results} is generated
mechanically from the \texttt{result.json} artifacts; no number in it is
transcribed by hand.

\subsection{Reference classes}
Every case is labelled by the kind of evidence it is checked against:

\begin{itemize}
\item[\textbf{[A]}] \textbf{Closed form or invariant.} An analytic solution, or a
  conservation law or objectivity property that must hold exactly. 21 cases.
\item[\textbf{[B]}] \textbf{Published numerical benchmark.} A tabulated result
  from the literature. 4 cases.
\item[\textbf{[E]}] \textbf{Experiment.} \emph{No case in this study is class
  [E].} See Section~\ref{sec:limits}.
\end{itemize}

\subsection{The tolerance protocol}
Three rules govern the study, and they exist because each has caught a real
error.

\begin{enumerate}
\item \textbf{Tolerances are declared before the solve.} The acceptance criterion
  is committed to version control before the case is run, and is never widened
  afterwards. A tolerance chosen after seeing the answer measures nothing.
\item \textbf{No unsourced reference values.} Every anchor is either derived
  symbolically in the case's own \texttt{reference.json} or cited to a specific
  table in a specific publication. Where no sourced value could be found, the
  case was reclassified downward rather than allowed to imply evidence it did not
  have --- this happened twice, both times from [E] to [A].
\item \textbf{Reported quantities must be mutually consistent.} The reference
  value, the computed value, and the error must refer to the same quantity.
  Section~\ref{sec:audit} reports where this rule was violated.
\end{enumerate}

\subsection{Verdicts}
A case is recorded as \textsc{pass}, \textsc{fail}, \textsc{partial}, or
\textsc{blocked}. \textsc{blocked} means the capability could not be assembled
from available input-file objects, or that the attempt exceeded its time budget;
it is distinct from \textsc{fail} and is never used to quietly retire a case that
did not work.

\subsection{A test that cannot fail is not evidence}
The governing methodological principle, stated separately because it drives the
design of most cases, is that agreement is only informative if disagreement was
possible. Several concrete failure modes recur: a probe placed where a boundary
condition pins the value it reports; a quantity measured over so short a segment
of a smooth curve that any nearby law would fit; a measurement quantised more
coarsely than its own tolerance; and a configuration whose symmetry fixes the
answer exactly, independent of the physics. Each of these has been found in this
study at least once, and one survived into the published results until the audit.

Where practical, cases are therefore designed so that accidental agreement is
implausible. The buckling case (Section~\ref{sec:buckling}) is the clearest
example.

\section{Results}
\label{sec:results}

All twenty-five cases meet their declared tolerance. Table~\ref{tab:main} reports
each case, the quantity scored, the measured error, the tolerance it was
measured against, and the fraction of the error budget consumed.

Consumption of the error budget is strongly skewed. Six cases consume less than
1\%, which reflects problems where the discretisation is essentially exact ---
the finite-strain objectivity check (T2-18) returns a rotation-stage stress
fourteen orders of magnitude below the elastic modulus, because an objective
formulation produces \emph{identically} zero stress under rigid rotation and the
residual is pure round-off. At the other end, four cases consume more than half
their budget and are marked $\dagger$ in Table~\ref{tab:main}. Those four are the
ones a reader should interrogate first; Section~\ref{sec:audit} does so.

{\footnotesize
\setlength{\tabcolsep}{4pt}
\begin{longtable}{@{}llcp{4.15cm}p{3.25cm}rrr@{}}
\caption{The twenty-five verified capabilities. ``Budget'' is the measured error
as a percentage of the declared tolerance; lower is better. $\dagger$ marks cases
consuming more than half their budget.\label{tab:main}}\\
\toprule
Case & Item & Cl. & Capability & Scored metric & Error & Tol. & Budget \%\\
\midrule
\endfirsthead
\toprule
Case & Item & Cl. & Capability & Scored metric & Error & Tol. & Budget \%\\
\midrule
\endhead
\bottomrule
\endfoot
\multicolumn{8}{@{}l}{\rule{0pt}{1.0em}\itshape 1.\ Solid mechanics and fracture}\\[1pt]
\texttt{T1-05} & 1.1 & A & J2 / von Mises elastoplasticity, thick cylinder & relative error in limit pressure & 0.00547 & 0.02 & 27.35\\
\texttt{T1-06} & 1.2 & A & Power-law secondary creep & relative error in final creep strain & $2.17\times 10^{-4}$ & 0.01 & 2.17\\
\texttt{T3-23} & 1.3 & A & Irradiation swelling eigenstrain and mismatch stress & L2 error in $\sigma_{xx}$, normalized & $7.28\times 10^{-5}$ & 0.001 & 7.28\\
\texttt{T1-08} & 1.4 & A & Hertzian elastic contact & relative error in p$_0$ and half-width a & $7.58\times 10^{-4}$ & 0.03 & 2.53\\
\texttt{T2-17} & 1.5 & B & Fracture mechanics, J and interaction integrals & relative J-integral error & 0.00362 & 0.03 & 12.07\\
\texttt{T2-18} & 1.6 & A & Finite-strain elastic kinematics & rotation-stage stress / elastic modulus & $1.88\times 10^{-14}$ & $1.00\times 10^{-8}$ & $1.9\times 10^{-4}$\\
\texttt{T3-22} & 1.7 & A & Euler buckling from an imperfect column & relative error in Southwell-recovered P$_{cr}$ & $1.77\times 10^{-4}$ & 0.03 & 0.588\\
\multicolumn{8}{@{}l}{\rule{0pt}{1.6em}\itshape 2.\ Heat transfer}\\[1pt]
\texttt{T1-01} & 2.1 & A & Transient heat conduction & max normalized L2 profile error, three times & $6.87\times 10^{-6}$ & 0.002 & 0.344\\
\texttt{T1-03} & 2.2 & A & Convective and radiative thermal boundary conditions & fin tip temperature and base heat rate & $3.37\times 10^{-4}$ & 0.01 & 3.37\\
\texttt{T2-21} & 2.3 & A & One-phase Stefan melting with apparent heat capacity & relative front-position error & 0.00287 & 0.02 & 14.34\\
\multicolumn{8}{@{}l}{\rule{0pt}{1.6em}\itshape 3.\ Fluid dynamics}\\[1pt]
\texttt{T1-07} & 3.1 & A & Laminar plane-channel flow & normalized L2 velocity-profile error & 0.00172 & 0.02 & 8.58\\
\texttt{T2-13} & 3.2 & B & Incompressible lid-driven cavity & L2 profile deviation / U$_{lid}$ & 0.0202 & 0.04 & 50.42\,$\dagger$\\
\texttt{T2-14} & 3.3 & B & Boussinesq natural convection & relative Nu and midplane-velocity error & 0.0349 & 0.04 & 87.32\,$\dagger$\\
\texttt{T2-15} & 3.4 & A & Conjugate heat transfer across a fluid/solid interface & relative Nusselt error & $1.03\times 10^{-4}$ & 0.02 & 0.515\\
\multicolumn{8}{@{}l}{\rule{0pt}{1.6em}\itshape 4.\ Porous flow}\\[1pt]
\texttt{T2-16} & 4.1 & A & Single-phase Darcy flow and Theis drawdown & relative drawdown error & 0.0111 & 0.02 & 55.51\,$\dagger$\\
\multicolumn{8}{@{}l}{\rule{0pt}{1.6em}\itshape 5.\ Phase field}\\[1pt]
\texttt{T1-09} & 5.1 & A & Split Cahn--Hilliard phase-field evolution & absolute error in the log--log slope & 0.0256 & 0.05 & 51.12\,$\dagger$\\
\texttt{T1-10} & 5.2 & A & Allen--Cahn curvature-driven grain growth & deviation of the R$^2$(t) fit from linearity & $3.72\times 10^{-5}$ & 0.001 & 3.72\\
\texttt{T3-25} & 5.3 & A & Anisotropic solidification, growth-direction selection & max error in arm direction, degrees & 0.168 & 2 & 8.40\\
\multicolumn{8}{@{}l}{\rule{0pt}{1.6em}\itshape 6.\ Mass transport}\\[1pt]
\texttt{T1-02} & 6.1 & A & Fickian species diffusion & max normalized L2 profile error, three times & $5.59\times 10^{-5}$ & 0.002 & 2.79\\
\texttt{T1-04} & 6.2 & A & Steady reaction--diffusion with a first-order sink & max normalized L2 profile error & $1.90\times 10^{-6}$ & 0.005 & 0.038\\
\multicolumn{8}{@{}l}{\rule{0pt}{1.6em}\itshape 7.\ Neutronics}\\[1pt]
\texttt{T1-11} & 7.1 & A & One-group neutron diffusion eigenvalue & relative k-effective error & $3.99\times 10^{-4}$ & 0.001 & 39.89\\
\multicolumn{8}{@{}l}{\rule{0pt}{1.6em}\itshape 8.\ Multiphysics coupling}\\[1pt]
\texttt{T2-20} & 8.1 & A & Neutronics-shaped heat source to conduction & relative profile L2 error & $1.85\times 10^{-10}$ & 0.01 & $1.8\times 10^{-6}$\\
\texttt{T3-24} & 8.2 & B & Fluid--structure interaction, Turek \& Hron FSI1 & relative tip displacement error & 0.00683 & 0.1 & 6.83\\
\multicolumn{8}{@{}l}{\rule{0pt}{1.6em}\itshape 9.\ Optimisation and uncertainty}\\[1pt]
\texttt{T2-19} & 9.1 & A & TAO optimization and inverse-parameter recovery & relative recovered-source error & $2.58\times 10^{-4}$ & 0.001 & 25.82\\
\texttt{T1-12} & 9.2 & A & Stochastic tools and Sobol sensitivity & absolute Sobol-index error & 0.0115 & 0.03 & 38.18\\
\end{longtable}
}


\section{Selected cases}

Four cases are described in more detail, chosen because each illustrates a
different way a verification case can be got wrong.

\subsection{Designing a test that cannot accidentally agree}
\label{sec:buckling}
The Euler buckling case (T3-22) recovers a critical load from an
imperfection-seeded nonlinear column by Southwell extrapolation~\cite{southwell1932}
and compares it against $P_{cr} = \pi^2 EI/(KL)^2$. A single boundary condition
would be weak evidence: a plausible-looking number could arise from a
mis-specified modulus or second moment. The case therefore runs two boundary
conditions whose effective-length factors differ by a factor of two, hence whose
critical loads differ by sixteen, and requires both to agree. It additionally
measures $EI$ independently by a transverse tip-load test, obtaining 666.508
against a nominal 666.667. The measured error against the closed form is
$1.77\times10^{-4}$, or 0.59\% of budget.

An earlier version of this case was wrong in a way worth recording. A
\texttt{Pressure} boundary condition defaults to acting on the displaced mesh,
which makes it a follower load; the column then becomes Beck's column, which has
no static bifurcation at all. The error did not move under mesh refinement, which
is the signature of a wrong model form rather than a discretisation error.

\subsection{Ensemble estimation of a stochastic exponent}
The Cahn--Hilliard case (T1-09) measures the Lifshitz--Slyozov--Wagner coarsening
exponent~\cite{lifshitz1961} from the decay of gradient
energy~\cite{cahn1958}, against a target of $-1/3$. The difficulty is
that a single realisation contains only about twenty domains at the measurement
scale, so one merge event moves the observable by percent. Widening the tolerance
to accommodate that noise would have destroyed the test's discriminating power.

Instead the case declares ten random seeds in advance and scores the ensemble.
All ten were run to completion and all ten were scored; no seed was dropped and
no optional stopping was applied. The ensemble fit gives $-0.3078$ and the mean
of the ten per-seed fits gives $-0.3045$, either side of a tolerance of $0.05$
about $-1/3$; the agreement of two independent estimators is stronger evidence
than either alone. Mass is conserved to $\sim10^{-11}$ per seed against a
$10^{-8}$ requirement.

This case also produced the clearest instance of a window rule satisfied for the
wrong reason. The measurement window opens when the characteristic domain size
exceeds ten interface widths, but that size is computed \emph{from} the gradient
energy; while the field is still uniform noise the computed size is enormous and
crosses the threshold from above. The tell was a coefficient of determination of
0.64 --- a scaling law does not fit that badly.

\subsection{Re-scoping an unreachable observable}
The anisotropic solidification case (T3-25) was originally declared against
dendrite tip velocity and its scaling with undercooling. That observable is not
reachable with this model~\cite{kobayashi1993} on available hardware: latent heat
saturates any domain that fits in memory, the melt recalesces, and the tip
stalls. A steady tip would require roughly seventeen times the elements, at days
per run.

Rather than report a tip velocity measured over a window where no steady tip
exists, the case was re-scoped to a different property of the same model that
\emph{is} exactly predicted: the crystal's growth direction must follow the
imposed anisotropy axis, $\theta_m = \theta_0$ modulo $360/m$. The observable is
the angular Fourier content of the solid region, $\int w\cos(m\theta)\,dA$ and its
sine counterpart, which recovers orientation and amplitude exactly and avoids the
node-straddling error of an interface-crossing measurement. Rotating the imposed
axis rotates the measured crystal one for one to within $0.17^{\circ}$.

This is a genuine reduction in scope, and it is recorded as such: the
tip-velocity half of the original capability is not covered and no steady-growth
claim is made.

\subsection{Checking the load before the pressure}
The Hertzian contact case (T1-08) compares peak contact pressure and contact
half-width against the Hertz solution. Two properties of penalty contact make a
naive comparison unsafe. First, nodal contact pressure carries a checkerboard
oscillation whose amplitude \emph{grows} under refinement, so a pointwise norm
against a smooth reference never converges even when the solution is correct.
Second, the applied load is quantised by the load-stepping scheme, so the
pressure inferred from a nominal load can be systematically off.

Both are handled by measuring rather than assuming: the traction integral is
asserted against the measured reaction before any comparison with the closed
form, and the Hertz half-width and peak pressure are evaluated at the load the
model actually developed rather than the load requested.

\section{Audit}
\label{sec:audit}

After the first twenty-four cases had been recorded as passing, the set was
audited against the three protocol rules. The audit examined tolerances,
reference provenance, and quantity consistency for all of them, and examined in
detail the ten cases the first pass flagged. The twenty-fifth case completed
afterwards and is described in Section~\ref{sec:fsi}.

\textbf{No verdict changed.} Every verifier was scoring the quantity it claimed
to score, against the tolerance it had declared. What the audit found was one
vacuous check and a systematic presentation defect.

\subsection{A check that could not have failed}
The mesh-independence check in the solidification case reported a phase spread of
$3.6\times10^{-15}$ across three meshes --- and this was not evidence of mesh
independence. The sweep had been run at an imposed axis of $90^{\circ}$, which is
$30^{\circ}$ modulo the $60^{\circ}$ period of the six-fold observable. At that
angle the configuration is mirror-symmetric about a mesh axis, the sine
coefficient vanishes identically, and the phase is exact \emph{by symmetry}
rather than by measurement.

The case's own reference file already recorded this symmetry property. It had
simply never been applied to the convergence block. Re-running the sweep at
$15^{\circ}$, which is not mesh-symmetric, produces a phase that
\emph{converges} on the imposed angle instead of being pinned to it:
$14.328^{\circ}$, $14.832^{\circ}$, $14.906^{\circ}$ at increasing resolution,
errors of $0.672^{\circ}$, $0.168^{\circ}$ and $0.094^{\circ}$. This is a check
that can fail, and does not.

The general lesson is narrower than ``check your symmetries'' and more useful: a
caveat recorded in one part of an artifact must be applied to every check in it.
The knowledge needed to catch this had been written down months earlier.

\subsection{Presentation defects}
Seven of the results reported a triple --- reference value, computed value, error
--- in which the error was not the difference of the first two. In four cases
this is legitimate and was simply not disclosed: the scored metric is a whole-field
profile norm, while the reference and computed columns carried a single
representative scalar. In the remaining three the artifact was wrong: two
reported a null reference beside a non-null error, and one placed a peak
temperature in the computed-value field while the error field held a profile
norm.

The most misleading instance was in the solidification case, where the computed
value stored raw principal values of an arctangent, so a true $30^{\circ}$
appeared as $-30^{\circ}$ and a true $45^{\circ}$ as $-14.855^{\circ}$. The
tolerance had always been declared modulo the observable's period, so the verdict
was correct --- but the summary table printed what read as two flat failures
beside a pass. All seven have been corrected, and the summary table now names the
scored metric and its tolerance explicitly.

\subsection{Caveats recorded but not corrected}
Three findings are properties of the results rather than of the reporting, and
are recorded rather than fixed.

\begin{itemize}
\item \textbf{Natural convection (T2-14) has not plateaued.} At $Ra=10^6$ the
  Nusselt error runs $16.11 \rightarrow 6.47 \rightarrow 4.32 \rightarrow 3.49\%$
  across four meshes against a 4\% tolerance. The sequence is monotone and
  heading toward the reference~\cite{devahldavis1983}, but the rate is close to
  first order rather than the second order the discretisation should deliver.
  The value 3.49\% is the error at $128^2$, not a converged value. This is the
  tightest margin in the study.
\item \textbf{The cavity case (T2-13) had its scoring rule changed after the
  fact.} The original rule scored the finest mesh and returned 0.1475, a
  failure. The current rule scores the most time-converged run and returns
  0.0202. Because the reference~\cite{ghia1982} is a steady solution, an
  under-converged fine mesh is genuinely further from it than a converged coarse
  one, so the change is physically defensible --- but the tolerance was met by
  changing what is measured, and the finest mesh remains unconverged.
\item \textbf{The coupling case (T2-20) converges faster than its own method.}
  Errors fall by a factor of sixteen per mesh doubling where four is expected.
  A rate that beats the method indicates that the manufactured source is nearly
  representable in the element space, so the case verifies the
  source-to-conduction mapping rather than the solver's spatial accuracy. It is
  valid for what it claims and should not be cited for more.
\end{itemize}

\subsection{Claims the audit cleared}
Two suspicions were checked and dismissed, and are reported because a null result
from an audit is also information. The grain-growth case (T1-10) scores a
linearity statistic, which raised the concern that any smooth curve would satisfy
it. Fitting the same statistic to synthetic data spanning the same range under
alternative growth exponents gives $4.4\times10^{-3}$ for a cubic law and
$1.1\times10^{-3}$ for an exponent of 2.5, both rejected by the $10^{-3}$
threshold that the true quadratic law clears at $2.5\times10^{-4}$; the window is
wide enough to discriminate. Separately, the inverse-problem case (T2-19)
recovers a source parameter from manufactured measurements, which invites the
charge of an inverse crime. The exact recovery visible at one mesh level is
indeed the mesh on which the measurements were generated, but the scored result
is taken on a finer mesh recovering from data generated elsewhere, and a noise
study is included.

\section{Threats to validity}
\label{sec:limits}

\textbf{This study contains no experimental validation.} Every case is checked
against a closed form or a published numerical benchmark. Two cases originally
carried an [E] label and were reclassified to [A] during the study, because in
each the reference was a self-declared illustrative figure rather than a sourced
experimental correlation. The correct summary of this work is therefore
``analytical and benchmark verification'', and any stronger claim would be
unsupported by what is here.

\textbf{Closed-form references constrain the problems that can be posed.} An
analytic solution generally exists only for simple geometry, linear or weakly
nonlinear material behaviour, and idealised boundary conditions. Agreement in
that regime does not extend automatically to the regimes a production model
occupies.

\textbf{Cases are two-dimensional and largely single-physics.} Each isolates one
solver path so that a failure localises rather than being guessed at. That is
deliberate, and it is also a limit: the study says little about the behaviour of
several capabilities exercised simultaneously, which is where a multiphysics
framework is most likely to disappoint.

\textbf{One framework version.} All results are for a single snapshot. Nothing
here speaks to behaviour before or after it.

\textbf{Benchmark references inherit their own uncertainty.} The three class-[B]
cases are compared against numbers that are themselves computed, not measured,
and carry the discretisation error of the codes that produced them.

\section{Fluid--structure interaction}
\label{sec:fsi}

The twenty-fifth case is the Turek and Hron FSI1 benchmark~\cite{turek2006}: an
elastic flag attached behind a rigid cylinder in channel flow at $Re=20$,
reaching a steady deflection. It is reported separately because it is the only
case in the study with two-way coupling, and because it was the only case
recorded as \textsc{blocked} for an extended period --- not because the
capability was absent, but because the reference had not been obtained from the
primary source and the geometry had not been built.

\subsection{Validate the halves before coupling them}
The benchmark supplies decoupled sub-problems with their own reference values,
and both were run on the FSI1 mesh before any coupling was attempted, so that a
coupled failure could be localised rather than guessed at. The fluid-only
problem (rigid flag, $Re=20$) converges monotonically to 0.45\% on drag across
1614, 6203 and 24692 elements; the solid-only problem (flag under gravity, no
fluid) converges to 1.11\% on tip deflection across 150, 570 and 2395 elements.

Neither sub-problem is scored as a result. Their reference values were declared
in advance but their \emph{tolerances} were not, and inventing a threshold after
seeing the answer is precisely what rule~1 forbids. They are therefore treated as
gates, judged on a criterion that does not require a threshold: the error must
fall monotonically under refinement, which a fortunate coarse mesh cannot
satisfy. One quantity fails that criterion and is reported rather than relied
upon --- the fluid-only lift runs 2.55\%, 0.15\%, 0.72\%, non-monotone because
lift here is a small difference of large surface tractions on a body only
2.5\,mm off the channel centreline.

\subsection{The coupled result}
The coupled model uses a monolithic arbitrary Lagrangian--Eulerian formulation
with pseudo-solid mesh motion and a penalty condition tying the fluid velocity to
the solid velocity across the interface. Table~\ref{tab:fsi} reports the mesh
study against Table~13 of the benchmark.

\begin{table}[h]\centering\footnotesize
\caption{FSI1 mesh study. Displacement is the scored quantity; the forces are an
independent corroboration.\label{tab:fsi}}
\begin{tabular}{@{}rrrrrr@{}}
\toprule
Elements & $u_x(A)$ err & $u_y(A)$ err & drag err & lift err & interface slip\\
\midrule
1\,764  & 1.28\% & 7.69\% & 0.52\% & 0.95\% & $3.51\times10^{-8}$\\
6\,773  & 0.83\% & 1.15\% & 0.20\% & 0.83\% & $3.56\times10^{-8}$\\
27\,087 & \textbf{0.36\%} & \textbf{0.68\%} & \textbf{0.12\%} & \textbf{0.92\%} & $3.57\times10^{-8}$\\
\bottomrule
\end{tabular}
\end{table}

The scored quantity is the worse of the two displacement components at the finest
level: 0.68\% against the 10\% tolerance declared before the case was built.
Error falls monotonically in both components.

Two features of this result carry more weight than the displacement agreement
itself. First, the forces corroborate independently. Drag and lift are not
scored, and they are not obtained the way the fluid-only case obtains them: at an
interface node the fluid momentum equation is closed by the penalty condition
rather than by a Dirichlet condition, so the saved momentum residual there is
missing the very term that enforces the coupling. Instead the load on the flag is
recovered from the reaction at its clamped root, which equilibrium makes equal
and opposite to the fluid load because FSI1 places no body force on the solid.
That construction was derived before the comparison. A tip displacement that
agreed while the forces did not would have indicated a coupling error rather than
a result.

Second, the principal risk was measured rather than argued away. Penalty coupling
is mesh dependent, and a refining mesh can begin to leak; the case therefore
reports the mean tangential velocity on the interface at every level. It is flat
at $3.5\times10^{-8}$ against a mean inlet velocity of $0.2$, so the interface
condition is not degrading as the mesh refines.

\subsection{What this case is not}
FSI1 is the steady member of the benchmark family. FSI2 and FSI3 are periodic and
are not attempted; nothing here speaks to time-accurate coupled dynamics. The
interface uses a penalty formulation rather than a mortar or Lagrange-multiplier
method, and the penalty value is not derived --- its adequacy is argued from the
measured slip, which is evidence but not proof.

\section{Conclusions}

Twenty-five solver capabilities of a multiphysics finite-element framework were
verified across nine physics domains against closed-form solutions and published
benchmarks, under tolerances committed before the solves. All twenty-five meet
their declared tolerance, with a median error consuming 7.3\% of the allowed
budget.

The more transferable result concerns method. A verification suite is not made
trustworthy by the number of passing cases; it is made trustworthy by the
possibility of failure, and that possibility has to be engineered case by case
and then re-examined. The audit reported here found a check that had been passing
for the wrong reason for months, and it found it not by re-deriving the physics
but by asking of each check what would have had to happen for it to fail. The
answer, in that one case, was ``nothing''.

\begin{thebibliography}{99}\footnotesize

\bibitem{cahn1958} J.~W. Cahn and J.~E. Hilliard,
``Free energy of a nonuniform system. I. Interfacial free energy,''
\emph{Journal of Chemical Physics}, vol.~28, no.~2, pp.~258--267, 1958.

\bibitem{devahldavis1983} G.~de~Vahl~Davis,
``Natural convection of air in a square cavity: a benchmark numerical solution,''
\emph{International Journal for Numerical Methods in Fluids}, vol.~3, no.~3,
pp.~249--264, 1983.

\bibitem{ghia1982} U.~Ghia, K.~N. Ghia, and C.~T. Shin,
``High-Re solutions for incompressible flow using the Navier--Stokes equations
and a multigrid method,'' \emph{Journal of Computational Physics}, vol.~48,
no.~3, pp.~387--411, 1982.

\bibitem{kobayashi1993} R.~Kobayashi,
``Modeling and numerical simulations of dendritic crystal growth,''
\emph{Physica D}, vol.~63, no.~3--4, pp.~410--423, 1993.

\bibitem{lifshitz1961} I.~M. Lifshitz and V.~V. Slyozov,
``The kinetics of precipitation from supersaturated solid solutions,''
\emph{Journal of Physics and Chemistry of Solids}, vol.~19, no.~1--2, pp.~35--50,
1961.

\bibitem{permann2020} C.~J. Permann \emph{et al.},
``MOOSE: Enabling massively parallel multiphysics simulation,''
\emph{SoftwareX}, vol.~11, 100430, 2020.

\bibitem{southwell1932} R.~V. Southwell,
``On the analysis of experimental observations in problems of elastic
stability,'' \emph{Proceedings of the Royal Society of London A}, vol.~135,
pp.~601--616, 1932.

\bibitem{theis1935} C.~V. Theis,
``The relation between the lowering of the piezometric surface and the rate and
duration of discharge of a well using ground-water storage,''
\emph{Transactions, American Geophysical Union}, vol.~16, pp.~519--524, 1935.

\bibitem{turek2006} S.~Turek and J.~Hron,
``Proposal for numerical benchmarking of fluid--structure interaction between an
elastic object and laminar incompressible flow,'' in \emph{Fluid--Structure
Interaction}, Lecture Notes in Computational Science and Engineering, vol.~53,
Springer, pp.~371--385, 2006.

\end{thebibliography}

\end{document}
